\section{重正化微扰展开的有限性}\label{sec:fin}

现在假定重正化常数 $Z$、$Z_3$、$\tilde{Z}_3$
取得使场 $(\ren{A})_{\mu}$、$\ren{c}$、$\ren{\cbar}$
的关联函数在 $D=4$ 处的奇异性按重正化耦合的所有阶消去。
本节的目的就是证明：此时涉及体场且流时间为正的关联函数也有限，
从而不需要进一步的重正化。

\subsection{重正化微扰论}

首先需要把微扰展开按重正化参数与场重新组织。
在这种形式下，展开由重正化场的作用量 $S_0+\Delta S$ 生成，
其中包含为消去费曼图在 $D=4$ 处的极点所需的抵消项 $\Delta S$。

具体地，把总作用量
\begin{equation}
  \Stot=S+\Sgf+\Sgh+\Sfl+\Sflgh
  \label{eq:7.1}
\end{equation}
中的裸参数与裸场用重正化的量表达，并令
\begin{align}
  S_0=\left.\Stot\right|_{Z=Z_3=\tilde{Z}_3=1},
  \label{eq:7.2}\\[2.5pt]
  \Delta S=\Stot-S_0+\Delta\Sbc,
  \label{eq:7.3}
\end{align}
其中
\begin{align}
   \Delta\Sbc=2\int\rmd^Dx\,\tr\bigl\{
   (Z^{1/2}Z_3^{1/2}-1)L_{\mu}(0,x)(\Ar)_{\mu}(x)
   \notag\\[2pt]
   +
   (\tilde{Z}_3Z^{1/2}Z_3^{1/2}-1)\dbar(0,x)\ren{c}(x)\bigr\}.
   \label{eq:7.4}
\end{align}
于是边界条件为
\begin{equation}
  \left.B_{\mu}\right|_{t=0}=(\ren{A})_{\mu},
  \qquad
  \left.d\right|_{t=0}=\ren{c},
  \label{eq:7.5}
\end{equation}
而由作用量 $S_0$ 导出的费曼规则因此与第 \ref{sec:perturbation}
节讨论的那些一致
（只是裸参数 $g_0$ 与 $\lambda_0$ 分别换成 $\mu^{\eps}g$ 与 $\lambda$）。

注意，边界条件 \eqref{eq:7.5} 与施加在裸场上的条件不同
（少了一个重正化因子的乘积）。为了纠正这一点，
必须在重正化场的作用量中计入抵消项 $\Delta\Sbc$。
这相当于在费曼规则中添加两点顶点，
它们对规范场及鬼场传播子的效果，
等价于让边界条件改变所缺失的那几个重正化因子。

\subsection{体抵消项的不存在}\label{ssec:nobulk}

在重正化微扰展开中，$D=4$ 处的奇异性（如果存在）
最早出现在某个圈阶 $l\geq1$。由于相互作用是定域的，
而且所有传播子在位置空间中都是缓增分布、其奇异性只出现在
坐标重合处，这一圈阶上关联函数的发散部分预计可以由定域抵消项消去。
与 Symanzik 研究过的半空间上普通场论的情形 \cite{Symanzik} 类似，
这些抵消项既可以定域在体内，也可以定域在半空间的边界上。

在本文考虑的理论中，体抵消项从一开始就可以排除。
为说明这一点，先注意：在大流时间处，
体场 $B_{\mu}$、$L_{\mu}$、$d$、$\dbar$ 的关联函数
由仅由流线与流顶点构成的费曼图给出。
所有这类图都是有向树或这类树的乘积，
每棵树终结于所考虑关联函数中某个 $B$ 或 $d$ 场上；
而且树另一端的流线必定起自关联函数中的 $L$ 或 $\dbar$ 场。

由于不存在圈图，体场的关联函数在大流时间处无奇性，
不需要重正化。因此发散的体抵消项被排除。
顺带指出：由 $B$ 场与 $d$ 场构成的定域复合场的关联函数同样有限，
因为当其中某些场的坐标重合时并不会生成圈图。
例如出现在 BRS 恒等式 \eqref{eq:6.23} 中的场 $D_{\mu}d$
就属于这种类型，因而无需重正化。

\subsection{边界抵消项}

上一小节的讨论表明：关联函数任何发散部分的结构
必然对应于一个定域在流时间零处的抵消项的插入。
既然重正化场 $(\ren{A})_{\mu}$、$\ren{c}$、$\ren{\cbar}$
的关联函数到所有阶都有限，发散只能来自至少含一个流顶点、
即至少一个带外流线的此类顶点的图。
因此可能的抵消项必正比于 $L_{\mu}$ 或 $\dbar$。

由于规范耦合无量纲，抵消项不得包含量纲大于 $4$
的复合场（那样的项是非相关的）。此外还必须尊重洛伦兹对称性、
整体规范不变性与鬼数守恒。
既然 $L_{\mu}$ 与 $\dbar$ 的量纲为 $3$、其他所有场的量纲为 $1$，
可知 $l$ 圈阶可能的抵消项形如
\begin{equation}
   2g^{2l}\int\rmd^Dx\,\tr\bigl\{
   z_1L_{\mu}(0,x)(\Ar)_{\mu}(x)+
   z_2\dbar(0,x)\ren{c}(x)\bigr\},
   \label{eq:7.6}
\end{equation}
其中系数 $z_1$、$z_2$ 是（奇异的）常数。原则上还应计入
一项 $LB$ 与一项 $\dbar d$；但鉴于边界条件 \eqref{eq:7.5}，
这些项并不独立，把它们包括进来等价于改动系数 $z_1$ 与 $z_2$。

\subsection{BRS 对称性的推论}

现在证明：该理论的 BRS 对称性把边界抵消项 \eqref{eq:7.6} 排除掉。
用重正化场与参数表述，BRS 恒等式 \eqref{eq:6.23} 读作
\begin{equation}
   \lambda\langle
   B^a_{\mu}(t,x)\partial_{\nu}(\ren{A})^b_{\nu}(y)
   \partial_{\rho}(\ren{A})^c_{\rho}(z)
   \rangle=
   -\langle(D_{\mu}d)^a(t,x)
   (\ren{\cbar})^b(y)\partial_{\rho}(\ren{A})^c_{\rho}(z)\rangle.
   \label{eq:7.7}
\end{equation}
鬼场 $c$、$\cbar$ 的非常规非对称重正化 \eqref{eq:5.13}
正是为了让重正化因子在此方程中不出现而选取的。
从 $\SUn$ 规范理论的角度看，鬼场的标准重正化与非对称重正化等价；
但在 $D+1$ 维理论中情形不同，因为只有 $c$ 与体场耦合。

式 \eqref{eq:7.7} 在尚需添加额外抵消项之前成立，
也就是直到圈阶 $l$ 为止（含）都成立。注意，
该方程各关联函数中的所有场——包括复合场 $D_{\mu}d$——
都是不能再被额外重正化的重正化场
（参见小节 \ref{ssec:nobulk}）。
因此，如果关联函数在 $l$ 圈阶奇异，
就必须能用形如 \eqref{eq:7.6} 的抵消项消去奇异性。

抵消项 \eqref{eq:7.6} 对关联函数的贡献，
通过把相应的两点顶点插入图~\ref{fig:brsdiag} 的树图得到。
这些插入的效果是使图 $1-5$ 的值分别乘上 $g^{2l}$ 倍的
\begin{equation}
   2z_1,\;z_1,\;z_1+z_2,\;z_1+z_2\;\hbox{以及}\;z_2.
   \label{eq:7.8}
\end{equation}
然而，回顾附录~\ref{app:C} 引用的各图的数值可以发现，
若 $z_1=z_2=0$ 不成立，以因子 \eqref{eq:7.8} 加权的各图之和
便违反 BRS 恒等式 \eqref{eq:7.7}。

因此式 \eqref{eq:7.7} 中的关联函数在 $l$ 圈阶必有限，
而所有其他（重正化的）关联函数也必无奇性，
因为向作用量添加形如 \eqref{eq:7.6} 的抵消项已被排除。
我们由此证明了：$D+1$ 维中的理论不需要进一步的重正化。
